%% The Elicit Author manual — UC-039, "Consult the Author's Manual".
%%
%% Build it with docs/manual/build-manual.sh, which renders it through a TeX Live container
%% and stamps it with the application version and build date (UC-039 BR-003). See README.md.

\documentclass[11pt, a4paper]{article}

\usepackage[T1]{fontenc}
\usepackage[a4paper, margin = 25mm, bottom = 28mm]{geometry}
\usepackage{elicit-brand}
\usepackage[hidelinks]{hyperref}

\hypersetup{
  pdftitle = {Elicit Author — Building a Survey},
  pdfauthor = {Elicit Software},
  pdfsubject = {Step-by-step instructions for building a survey in Elicit Author},
  colorlinks = true,
  linkcolor = elicitInkDark,
  urlcolor = elicitInkDark,
}

\graphicspath{{./}}

\begin{document}

%% ------------------------------------------------------------------ title page
\begin{titlepage}
\thispagestyle{empty}
\centering
\vspace*{3.5cm}

\includegraphics[width = 34mm]{images/elicit-logo.png}

\vspace{1.4cm}

{\color{elicitInkDarkest}\Huge\bfseries Building a Survey\par}

\vspace{0.5cm}

{\color{elicitSlate}\Large The Elicit Author manual\par}

\vspace{2.6cm}

{\color{elicitSlate}\large
Version \textbf{\elicitVersion}\\[0.35em]
Built \textbf{\elicitBuildDate}\par}

\vfill

{\color{elicitSlate}\small
This manual was produced by the same build as the application it ships inside,\\
and describes that version. Elicit Software.\par}

\vspace{1cm}
\end{titlepage}

%% ------------------------------------------------------------------ contents
\pagenumbering{roman}
\setcounter{page}{1}
{\color{elicitInkDarkest}\section*{Contents}}
\vspace{-0.6em}
\begingroup
\hypersetup{linkcolor = elicitCharcoal}
\makeatletter\@starttoc{toc}\makeatother
\endgroup

\clearpage
\pagenumbering{arabic}
\setcounter{page}{1}

%% ------------------------------------------------------------------ 1
\section{What a survey is made of}
\markboth{What a survey is made of}{}

An Elicit survey is a decision tree a respondent walks through. What they answer decides what
they see next, so designing a survey means designing that path. Everything in this manual uses
the same names and the same glyphs the application puts on screen.

\subsection{The containers}

Everything sits inside something else:

\begin{itemize}
  \item A \term{survey} holds steps.
  \item A \term{step} is a page in the respondent's progress bar. On the designer board it is a
        lane that groups related sections.
  \item A \term{section} is a titled group of questions on that page. On the board it is a card
        inside a lane.
  \item A \term{question} is one input the respondent answers. On the board it is a row inside a
        card, carrying a glyph that shows its type.
\end{itemize}

Steps run in order, sections run in order within their step, and questions run in order within
their section. Unless a rule says otherwise, that display order is exactly what the respondent
experiences.

\subsection{The rules}

A \term{rule} reads the answer to one upstream question, compares it with a reference value using
an operator, and when the comparison holds performs an action on one downstream target — a step,
a section or a question. Every rule is the same sentence:

\begin{quote}
\ui{question} \; (\ui{operator}) \; \ui{value} \; $\rightarrow$ \; (\ui{action}) \; \ui{target}
\end{quote}

There are three actions:

\begin{itemize}
  \item \term{Show} — the target appears only when the rule holds. Until then the respondent
        never sees it.
  \item \term{Repeat} — the target section or question is repeated as many times as the answer
        says, for example one \ui{Child} section per child.
  \item \term{Text} — the answer is written into the target's text wherever a
        \texttt{\{<TOKEN>|default\}} placeholder appears, so a later question can say
        ``your daughter'' rather than ``your child''.
\end{itemize}

There is no Hide rule. Hiding is the same as not showing: give the element a Show rule and it
stays hidden until that rule holds. The operators are Boolean, Greater than, Equal, Not equal,
Field exists and Contains. When several rules point at the same target, all of them must hold.

\begin{elicitnote}[The same page, inside the application]
Everything in this section is also on the \ui{How surveys work} page in the application
navigation, with the glyphs drawn live. This manual is its printable form, and adds the
step-by-step procedure that follows.
\end{elicitnote}

\subsection{Identity and publishing}

Every element gets a key the moment it is created and keeps it for life. Removing an element
\term{retires} it rather than deleting it, so answers already recorded at a site keep their
meaning. When the survey is ready you \term{export} a \texttt{.elicit} file; the Admin app
installs it at a site the first time and updates it in place after that, matching elements by
key.

\begin{elicitnote}[Vocabulary]
Author never uses ``publish'' for a state in the database — the verb is \emph{export}. The
credential a respondent enters to reach a survey is the \term{access code}, never a token;
``token'' means only the \texttt{\{<KEY>|default\}} placeholder in question text.
\end{elicitnote}

%% ------------------------------------------------------------------ 2
\section{Signing in}
\markboth{Signing in}{}

Author uses the platform's identity provider, the same one as the Admin and Survey apps. Open
Author in a browser and you are sent to the sign-in page.

\screenshot{01-sign-in}{The identity provider's sign-in page. Author itself never asks for a
password.}

Authoring requires the \texttt{elicit\_author} role. It is granted explicitly and is \emph{not}
implied by \texttt{elicit\_admin}: an administrator who has not been given it can open the
System pages but not a survey. If you sign in and the surveys are not there, that role is the
first thing to check with whoever administers your identity provider.

After signing in you land on the \ui{Surveys} list, the application's front door.

\screenshot{02-surveys-list}{The Surveys list. Every survey you can edit, with who last edited
it and the note from its last export.}

%% ------------------------------------------------------------------ 3
\section{Creating a survey}
\markboth{Creating a survey}{}

Choose \ui{New survey}. A survey needs two things to exist: a \ui{Name}, which is the machine
name and may not contain spaces, and a \ui{Title}, which is what a respondent sees.

\screenshot{03-new-survey}{Creating a survey. Name and Title are required; the rest can be filled
in later from \ui{Edit details}.}

\begin{itemize}
  \item \ui{Name} — the machine name, for example \texttt{FamilyHistory}. No spaces.
  \item \ui{Title} — the respondent-facing title, for example ``Family Health History''.
  \item \ui{Description} — a note for your colleagues; respondents never see it.
  \item \ui{Initial display key} — where the respondent starts. Leave it empty unless you have a
        reason not to.
  \item \ui{Post-survey address} — where the respondent is sent when they finish.
\end{itemize}

Save, and the survey appears in the list. Open it and you get its overview: the survey key, when
it was last exported, the four things you can do to it, and a validation panel.

\screenshot{04-survey-overview}{A survey's overview page. The validation panel is the thing to
watch — it is what stands between you and a clean export.}

The validation panel is live. It tells you what is still wrong with the definition — a survey
with no steps, a choice question with no list, a rule that points backwards — and it has to read
``Ready to export: no findings'' before the export is worth doing.

%% ------------------------------------------------------------------ 4
\section{Laying out steps and sections}
\markboth{Laying out steps and sections}{}

Choose \ui{Design…} from the survey overview. The drawer on the left is your toolbox; the board
on the right is the survey; the panel beside the board previews the selection as a respondent
will see it.

\screenshot{05-designer-board}{The designer. Controls, Lists, Sections, Questions and the
structure tree on the left; the board in the middle; the respondent preview on the right.}

Work top-down: steps, then sections, then questions, then the lists, rules and reporting tags
that give them meaning.

\subsection{Adding a step}

Use \ui{Add step} at the top of the board, or drag the \ui{STEP} glyph from \ui{Controls} onto
the board. Both open the same editor.

\screenshot{06-new-step}{A new step. \ui{Dimension name} is what the reporting schema calls this
step; it defaults to the name, and changing it later is a reporting change rather than a rename.}

\begin{elicitnote}[Everything a drag can do, a menu can do too]
Every element on the board carries a \ui{$\vdots$} menu with the same actions the drag gestures
perform, so the designer is usable without a pointer.
\end{elicitnote}

\screenshot{07-step-menu}{A step's menu: edit it, add another step after it, add a section, move
it up or down, open its details page, or remove it.}

\subsection{Adding a section}

A step with no sections shows the respondent nothing. Add one from the step's menu, or drag the
\ui{SECTION} glyph into the lane.

\screenshot{08-new-section}{A new section. The name is the page heading the respondent reads.}

One click on a lane opens it to reveal its section cards; one click on a card opens it to reveal
its question rows. A double click on any of them opens its editor.

%% ------------------------------------------------------------------ 5
\section{Giving choice questions their options}
\markboth{Giving choice questions their options}{}

A choice question — radio, combo box, multi-select or checkbox group — needs a \term{list of
options} before it can be saved. Define the list first, from the \ui{Lists} panel in the drawer.

\screenshot{09-list-of-options}{A new list of options. The options are in the order respondents
see them; the arrows reorder them and the coded value defaults to the display text.}

A list belongs to the survey, not to one question, so several questions can offer the same
options. Drag a list's glyph onto a question to give it those options, or pick the list in the
question's own editor.

%% ------------------------------------------------------------------ 6
\section{Adding questions}
\markboth{Adding questions}{}

Add a question from a section's menu, or drag the \ui{QUESTION} glyph — set to the type you want
— into the card.

\screenshot{10-question-editor}{A text question. \ui{Short text} is what the board and the
structure tree show; \ui{Question text} is what the respondent reads.}

The fields worth knowing:

\begin{itemize}
  \item \ui{Type} — text, integer, radio, combo box, date and the rest. The editor changes with
        the type, and the preview at the bottom shows the field the respondent will get.
  \item \ui{Short text} — the label on the board and in the structure tree. Keep it short.
  \item \ui{Question text} — what the respondent reads. This is where a
        \texttt{\{<TOKEN>|default\}} placeholder goes if a Text rule is to fill it.
  \item \ui{Required} and \ui{Validation message} — whether an answer is compulsory, and what
        the respondent is told when it is missing.
  \item \ui{Sample answer} — previews what a Text rule reading this question would write
        elsewhere.
\end{itemize}

Choose a choice type and the editor asks for the list of options you defined earlier.

\screenshot{11-choice-question}{A radio question with its list of options attached. Without a
list the editor refuses to save.}

Once a section is open you can see the whole page taking shape: rows in order, each with its type
glyph, and the respondent's view of the same page beside it.

\screenshot{12-section-opened}{A section opened on the board, with the rules drawn as arrows and
the respondent preview beside it.}

\subsection{Asking the same question twice}

The \ui{Sections} and \ui{Questions} panels in the drawer list every section and every question
in the survey once, with the number of places each appears in. Drag one into another container,
or use \ui{Also add to section…}, and it is the \emph{same} element shown in a second place — a
change to its text reaches every section it is in.

\screenshot{15-question-bank}{The Questions panel. Each question appears once, with how many
sections use it; the search box narrows the panel as you type.}

%% ------------------------------------------------------------------ 7
\section{Connecting questions with rules}
\markboth{Connecting questions with rules}{}

Anything conditional is a rule. Start one from the upstream question — the one whose answer is
read — using \ui{Add rule…} in its menu, or by dragging a rule glyph from \ui{Controls} onto the
question and then dragging the arrowhead onto the target.

\screenshot{13-question-menu}{A question's menu. Rules and reporting tags both start here.}

The editor is the rule sentence, in the order the runtime evaluates it.

\screenshot{14-new-rule}{A Repeat rule: when \ui{Child count} is greater than 0, repeat the
section \ui{Child} that many times.}

\begin{elicitcaution}[Rules only point forwards]
The target must come after the upstream question in the survey's order. Author refuses a rule
that points backwards and says so, because the respondent would reach the target before there is
an answer to read. A count question therefore sits in the section \emph{before} the block it
repeats — a section cannot repeat itself.
\end{elicitcaution}

Rules are drawn on the board as arrows from the question that is read to the element that is
acted on, labeled with the action. Untick \ui{Show rules} at the top of the board to hide them
while you work on structure.

%% ------------------------------------------------------------------ 8
\section{Tagging for reporting}
\markboth{Tagging for reporting}{}

A \term{reporting tag} names a column of the report a site builds from the answers. Tag a
question and its answers fill that column; tag a section as it appears in a step with a fixed
value, and every respondent who reaches it is classified by that value. Tags that share a
\term{dimension} share one table of values.

Open the survey's \ui{Reporting…} page to see every tag, every assignment, and the schema they
will produce.

\screenshot{16-reporting}{The reporting page: the tags, where they are assigned, and the
dimension tables and report columns the site's reporting build will create.}

Define a tag with \ui{New tag…} and attach it with \ui{Assign tag…}, or use \ui{Reporting tags…}
in a question's or section's own menu.

\screenshot{17-assign-tag}{Assigning a tag. A question can report the respondent's answer; a
section in a step can classify everyone who reaches it by a fixed value.}

\begin{elicitcaution}[Tag changes are retroactive]
Changing or removing an assigned tag relabels answers that deployed sites have \emph{already}
recorded, once the survey is exported and applied. Historical reports change with it. Treat a
tag change as a decision about past data, not only about future data.
\end{elicitcaution}

%% ------------------------------------------------------------------ 9
\section{Checking and exporting}
\markboth{Checking and exporting}{}

Go back to the survey overview and read the validation panel. It must say ``Ready to export: no
findings''. Anything else names the element and the problem — fix it on the board and the panel
updates.

Then choose \ui{Export…}.

\screenshot{18-export}{Exporting. The release note is required: it is recorded on the survey and
inside the file, and it is what a colleague sees in the Surveys list.}

The export writes a \texttt{.elicit} file in the \texttt{ELICIT\_SURVEY\_EXPORT\_V1} format,
records who exported it and the note, and mints a revision. Hand that file to the Admin app's
\ui{Apply Survey Definition}: the first time it installs the survey at a site, and every time
after that it updates it in place, matching elements by their keys. Elements you removed travel
as retirements, so answers already recorded keep their meaning.

\begin{elicitnote}[The round trip]
\ui{Import definition} in the application navigation reads a \texttt{.elicit} file back into
Author for further editing, keeping every key. That is how you pick up a survey someone else
exported, or start from one of the sample definitions.
\end{elicitnote}

%% ------------------------------------------------------------------ 10
\section{Where to look next}
\markboth{Where to look next}{}

\begin{itemize}
  \item \ui{How surveys work} in the application navigation — this manual's first section, drawn
        live with the real glyphs.
  \item \ui{Import definition} — bring an existing \texttt{.elicit} file into Author.
  \item \ui{System} (for administrators) — version, database, branding, connections and OIDC
        diagnostics for a deployment.
\end{itemize}

\screenshot{19-guide}{The in-app guide. It is always the version you are running, because it
ships with the application — as does this manual.}

\vspace{1em}

\begin{elicitnote}[Which version is this?]
This manual describes Elicit Author \textbf{\elicitVersion}, built \textbf{\elicitBuildDate}.
The same pair is printed at the foot of every page, so a printed copy can always be matched back
to the build it came from.
\end{elicitnote}

\end{document}
